\section{Proofs for qualification-optimal readouts and sequential costs}\label{supp:proofs73}
\subsection{The readout program and its dual}\label{supp:readout73}
\begin{proof}[Proof of Theorem~\ref{thm:readout73}]
The readout has Lipschitz constant at most $L(\theta)=\sum_kL_k|\theta_k|$. At a cover point the residual is in $[y_j^--F_j\theta,y_j^+-F_j\theta]$. The Lipschitz constant of that residual is at most $C+L(\theta)$. For any state choose a cover point at distance at most $\rho$. Its lower endpoint minus $\rho(C+L(\theta))$ is a valid lower bound there, and the upper endpoint plus that quantity is a valid upper bound. Taking the minimum and maximum over the cover gives the asserted oscillation bound. No distribution of states, interpolation of a discontinuous actor, or convexity of the fitted continuation is used.

Introduce $u,l\in\mathbb R$ and $s_k\ge0$ and minimize
\[
 u-l+2\rho C+\sum_k\lambda_ks_k
\]
subject to
\begin{align*}
 u&\ge y_j^+-F_j\theta,& l&\le y_j^--F_j\theta &&(j=1,\ldots,N),\\
 s_k&\ge\theta_k,&s_k&\ge-\theta_k &&(k=1,\ldots,K),
\end{align*}
and $-B_k\le\theta_k\le B_k$. The stated dimensions follow. For any $\theta$ the tight choices of $u,l$ and the choice $s_k=|\theta_k|$ attain \eqref{eq:readout-objective73}; conversely every feasible point has objective at least that expression. Continuity on the compact coefficient box gives attainment.

The maximum of finitely many real numbers is the maximum of their averages over probability vectors. Consequently the objective can be written as
\[
 2\rho C+\max_{p,q\in\Delta_N}
 \left\{p^\top y^+-q^\top y^-+
       [F^\top(q-p)]^\top\theta+\sum_k\lambda_k|\theta_k|\right\}.
\]
For a scalar $z$,
\[
 \min_{|v|\le B}\{zv+\lambda|v|\}
       =-B(|z|-\lambda)_+.
\]
Indeed zero is minimizing when $|z|\le\lambda$, whereas the appropriate endpoint of the coefficient interval is minimizing otherwise. Applying this identity coordinatewise proves weak duality in \eqref{eq:readout-dual73}. The finite program is feasible and has a finite attained optimum. Linear-programming duality, or equivalently minimax on the compact coefficient box and the two probability simplexes, interchanges the minimum and maximum in the last display. It gives equality for some $p,q$. For any explicitly supplied feasible triple, subtracting its dual lower bound from its primal objective bounds its excess over the optimum.
\end{proof}

\paragraph{Rational verification and the backward qualification chain.}
The empirical implementation uses rational feature values, rational label endpoints and rational moduli. A numerical optimizer proposes coefficients and dual multipliers. Coefficients are rounded to a fixed dyadic lattice and clipped to the declared box. Nonnegative dual proposals are rounded and normalized to sum exactly to one; a fixed unit vector is the valid fallback if their mass vanishes. The objective, support function and optimization gap are then recomputed in rational arithmetic. Solver failure or a gap exceeding the declared tolerance is recorded rather than converted into an assertion of optimal fitting.

At date $t$ this argument conditions on the actual $h_{t+1}$ produced by the next date's computation. Repeating backwards gives valid bounds $\omega_t\le\mathcal W_t$ simultaneously because they are deterministic qualifications, not confidence statements. Substituting them into the frozen-continuation proof gives
\[
 \Omega_h\le\sum_{t<T}b_t\mathcal W_t+b_T\omega_T,
 \qquad J_0^\pi-V_0^*\le\sum_{t<T}b_t\mathcal W_t+b_T\omega_T+H_h^\pi.
\]
The same substitution is valid in the cost-contrast support. This does not compare the future continuations generated by different backward optimization choices, and hence makes no joint optimality assertion.

\subsection{Whole-cell policy qualification in the executed construction}\label{supp:construction73}
For the new two-state construction, the original terminal payoff is retained exactly. Each subsequent continuation is additive with nine coefficients: a constant and four features per state coordinate. The ReLU family uses $x_j$ and $(x_j-k/4)_+$, $k=1,2,3$. The conventional family uses the four nonconstant degree-four Bernstein terms $\binom4k x_j^k(1-x_j)^{4-k}$, $k=1,\ldots,4$. Feature Lipschitz bounds are respectively one and four. The latter follows by differentiating a Bernstein term: its derivative is four times the difference of two nonnegative degree-three Bernstein terms whose sum over the full basis is one. Both families have the same coefficient box and node information; their spanning spaces and modulus bounds need not coincide.

The midpoint cover with $n$ cells per coordinate has radius $\rho=1/n$ in $\ell^1$. If the actual future has Lipschitz constant $L_+$, the inherited primitive inequalities give
\[
 L_x=4+\beta\frac{11}{16}L_+,
 \qquad L_a=\frac{13}{16}+\beta\frac34L_+,
 \qquad C=L_x+\frac{L_a}{16},\quad\beta=\frac{15}{16}.
\]
At the terminal continuation $L_+=5$; at an earlier date it is the verified coefficient-weighted feature modulus. The 33-point action mesh at a midpoint $x$ has exact fill radius $c(x)/64$. Thus subtracting $L_ac(x)/64$ from its minimum action value gives the lower continuous-action label, while the minimum itself is the upper label. All those evaluations are rational conditional expectations under the original uniform innovation law.

For the implemented actor, a second 33-point menu is rounded down to the $1/4096$ lattice and bounded by the minimum capacity over the whole observation cell. Let $a_j$ minimize its exact midpoint action value and let $y_j^-$ be the continuous-action lower label. At any state $x$ in that cell,
\begin{align*}
 Q^h(x,a_j)-\inf_a Q^h(x,a)
 &\le Q^h(x_j,a_j)-y_j^-+(L_x+C)\|x-x_j\|_1\\
 &\le Q^h(x_j,a_j)-y_j^-+(L_x+C)\rho.
\end{align*}
Taking the maximum over all cells gives the recorded uniform greedy-gap allowance. Replacing $a_j$ by zero gives the comparator allowance. The selected action is feasible throughout its cell, including its closed boundary, and a fixed boundary convention gives a well-defined actor. This proves the continuous-state, continuous-action qualification; midpoint evaluations alone would not do so.

For a ReLU future, exact conditional integration follows Proposition~\ref{prop:moment72}. The independent checker instead splits the innovation interval at every hinge crossing and integrates each affine piece by its endpoint trapezoid. For a Bernstein future, the production calculation expands the uniform moments. The independent checker uses the positive five-point weights $(7,32,12,32,7)/90$ on equally spaced points from $-1/32$ to $1/32$. Substitution of each monomial of degree zero through five verifies this rule directly, so it integrates the conditional degree-four polynomial exactly. The checker separately uses the retained original rational terminal primitive at every final-date action value. These checks concern all new construction nodes and actions, not the entire historical recursive Bellman computation.

\subsection{A time-uniform endpoint bound}\label{supp:anytime73}
\begin{proof}[Proof of Theorem~\ref{thm:anytime73}]
Condition on $\mathcal F_{i-1}$. The frozen-continuation telescoping proof applies to the two complete controllers selected for this episode; it gives $\mathbb E[D_i\mid\mathcal F_{i-1}]=\Delta_i$. History-dependent actions cause no change in the tower property provided the actual conditional transition kernels and all required greedy-gap qualifications remain valid. The common initial law cancels the initial continuation term. The terminal correction is still required.

Write $R_i=b_i-a_i$. Since $L_i\le D_i\le U_i$ and the endpoints lie in $[a_i,b_i]$,
\[
 \mathbb E[L_i\mid\mathcal F_{i-1}]\le\Delta_i
       \le\mathbb E[U_i\mid\mathcal F_{i-1}].
\]
Conditional Hoeffding's lemma therefore gives, for each fixed $\lambda>0$,
\[
 \mathbb E\!\left[
   \exp\!\left\{\lambda(L_i-\Delta_i)-\frac{\lambda^2R_i^2}{8}\right\}
                    \middle|\mathcal F_{i-1}\right]\le1.
\]
To see the constant directly, a variable supported on an interval of length $R$ has tilted variance at most $R^2/4$. The second derivative of its centered log moment-generating function is this tilted variance. Integrating twice from zero bounds that log moment-generating function by $\lambda^2R^2/8$. The nonpositive conditional mean of $L_i-\Delta_i$ can only decrease the exponential bound. The same reasoning applies to $\Delta_i-U_i$.

Products of these conditional factors are nonnegative supermartingales starting at one. Ville's inequality bounds the probability that the lower-endpoint product for $\lambda_j$ ever exceeds $2/(\alpha w_j)$ by $\alpha w_j/2$. The same bound holds for the upper-endpoint product. A union bound over both signs and all $j$ gives simultaneous validity with failure probability at most $\alpha\sum_jw_j\le\alpha$. Outside that exceptional event every partial sum is bounded by every corresponding linear boundary. Taking their minimum gives \eqref{eq:anytime-boundary73} and \eqref{eq:anytime-interval73}. The predictable support also bounds each $\Delta_i$, justifying the additional intersection.

The event already covers every finite index, so evaluating it at a finite stopping time does not consume additional error probability. With a fixed complete controller pair and fixed initial-law/reset contract, its conditional expected original-law contrast is the same number $\Delta$ at every episode. Predictable changes of a valid critic alter the observations and ranges, not this identity. Without that fixed target, the preceding calculation identifies only the displayed conditional average.
\end{proof}

The implementation need not evaluate exponentials. For $x\ge1$, write $x=2^ky$ with $1\le y<2$, and use
\[
 \log y=2\sum_{r=0}^{m-1}\frac{z^{2r+1}}{2r+1}+\mathcal R_m,
 \qquad 0\le\mathcal R_m\le
       \frac{2z^{2m+1}}{(2m+1)(1-z^2)},\quad z=\frac{y-1}{y+1}.
\]
Use the same expression at $z=1/3$ for $\log2$. It gives an exact rational upper bound on each logarithm in the boundary; replacing the logarithm by that bound enlarges the interval. Arithmetic operation counts do not bound rational bit complexity.

The finite set of bets is deliberate: a finite minimum is directly computable. It remains time-uniform for arbitrarily many episodes, although a finite bet set can leave a positive asymptotic radius. Increasing the predeclared set improves its supported scales but is a new design choice, not a free after-the-fact selection. One may instead preallocate a countable summable family; an implementation using any finite subset retains validity with its assigned weights.

\subsection{Stopping, failed episodes and the replacement decision}\label{supp:adoption73}
\begin{proof}[Proof of Corollary~\ref{cor:adoption73}]
On the simultaneous event, $l_N\le\Delta\le u_N$. Adding the declared incremental charge and the one-time charge divided by the deployment count gives the claimed interval. Its upper endpoint being negative implies strict net cost reduction. Inclusion in an equivalence band is a different two-sided assertion. The event supports every such deterministic algebraic implication, including a stopping rule that inspects the interval repeatedly.

For a new epoch, condition on the history at its start, which includes the construction and selection of its new controller pair. Apply the theorem with that epoch's previsible error allocation $\alpha_e$. A conditional union bound followed by expectation bounds the probability of any erroneous qualification by $\mathbb E\sum_e\alpha_e$, no larger than the total declared budget. No independence between epochs is required.
\end{proof}

A support reserved before an episode remains valid when its numerical evaluation fails. Recording the full interval $[a_i,b_i]$ then preserves endpoint containment and the martingale argument. Deleting the episode after an outcome-dependent failure need not preserve the conditional law. All attempted work is retained irrespective of which interval is recorded. An invalid support certificate is different: it supplies no admissible substitute observation and must not be treated as a harmless numerical timeout.

Finally, a conditional kernel-bias bound can be accommodated only when it is itself justified: if $|\mathbb E[D_i\mid\mathcal F_{i-1}]-\Delta_i|\le\zeta_i$ with previsible certified $\zeta_i$, enlarge both ends by $n^{-1}\sum_{i\le n}\zeta_i$. This follows by applying the proved bound to the actual conditional mean and then using the triangle inequality. An empirical model discrepancy without such a bound is not set to zero.
